\subsection{Tests, Likelihood Ratios, and Error Control}

Testing formalises decisions between hypotheses.  The central quantities are the probability of false rejection under the null and the probability of rejection under alternatives.

\begin{definition} \label{def:hypotheses-tests-size-power} \hlt{Hypotheses, Tests, Size, and Power}
\begin{enumerate}[label=(\roman*)]
\setlength{\itemsep}{0pt}
\item A \hlt{null hypothesis} is a subset $\Theta_0\subseteq\Theta$ of the parameter space.
\item An \hlt{alternative} is a subset $\Theta_1\subseteq\Theta$ disjoint from $\Theta_0$, so $\Theta_0\cap\Theta_1=\emptyset$.
\item Randomised \hlt{test} is a measurable map $\varphi:\mathcal{X}\to[0,1]$, where $\varphi(x)$ is rejection probability at observation $x$.
\item The \hlt{size} of $\varphi$ is $\alpha(\varphi)=\sup\limits_{\theta\in\Theta_0}\mathbb{E}_\theta[\varphi(X)]$.
\item The \hlt{power function} of $\varphi$ is $\beta_\varphi(\theta)=\mathbb{E}_\theta[\varphi(X)]$.
\end{enumerate}
\end{definition}

\begin{definition} \label{def:simple-composite-type-errors} \hlt{Simple and Composite Hypotheses; Type Errors}
\begin{enumerate}[label=(\roman*)]
\setlength{\itemsep}{0pt}
\item A \hlt{simple} hypothesis specifies one law; a \hlt{composite} hypothesis specifies several.
\item A \hlt{type I error} rejects a true $H_0$; its probability is bounded by $\alpha(\varphi)$.
\item For a simple alternative $P_1$, a \hlt{type II error} is failure to reject $H_0$ under $P_1$, with probability $\mathbb{E}_1[1-\varphi(X)]$.
\item The \hlt{power} against $P_1$ is $1-\mathbb{E}_1[1-\varphi(X)]$.
\end{enumerate}
\end{definition}

\begin{proposition} \label{prop:posterior-odds-likelihood-ratio} \hlt{Posterior Odds and Likelihood Ratios}\\
Let $H_0,H_1$ be simple hypotheses with $\mathbb{P}(H_i)=\pi_i>0$ and likelihood $p_i(x)=p(x\mid H_i)$ for $i=0,1$.\\
For $m(x)\defeq\pi_0p_0(x)+\pi_1p_1(x)>0$, $\mathbb{P}(H_i\mid x)=\frac{\pi_ip_i(x)}{m(x)}$ and $\frac{\mathbb{P}(H_0\mid x)}{\mathbb{P}(H_1\mid x)}=(\frac{\pi_0}{\pi_1})(\frac{p_0(x)}{p_1(x)})$.
\begin{proof}
Bayes' rule gives $\mathbb{P}(H_i\mid x)=\frac{\pi_ip_i(x)}{m(x)}$ for $i=0,1$. Therefore $\frac{\mathbb{P}(H_0\mid x)}{\mathbb{P}(H_1\mid x)}=\frac{\frac{\pi_0p_0(x)}{m(x)}}{\frac{\pi_1p_1(x)}{m(x)}}=(\frac{\pi_0}{\pi_1})(\frac{p_0(x)}{p_1(x)})$.\\
\hlt{Structural Commentary}: Posterior odds factor as prior odds times the data likelihood ratio.\\
\hlt{Key Insight}: Bayes updating multiplies odds by $\frac{p_0(x)}{p_1(x)}$.
\end{proof}
\end{proposition}

\begin{theorem} \label{thm:neyman-pearson-lemma} \hlt{Neyman--Pearson Lemma}\\
Let $P_1\ll P_0$, $\Lambda\defeq \frac{dP_1}{dP_0}$, $\alpha\in[0,1]$. If $c\geq0$ and $\gamma\in[0,1]$ make $\varphi^\ast\defeq\mathbf{1}_{\{\Lambda>c\}}+\gamma\mathbf{1}_{\{\Lambda=c\}}$ satisfy $\mathbb{E}_0\varphi^\ast=\alpha$,\\
then $\varphi^\ast$ is most powerful among level-$\alpha$ tests: for every $\psi:\mathcal{X}\to[0,1]$ with $\mathbb{E}_0\psi\leq\alpha$, $\mathbb{E}_1\psi\leq\mathbb{E}_1\varphi^\ast$.
\begin{proof}
For any level-$\alpha$ competitor $\psi$, $(\varphi^\ast-\psi)(\Lambda-c)\geq0$ pointwise.\\
Since $dP_1=\Lambda\,dP_0$, $\mathbb{E}_1(\varphi^\ast-\psi)=\mathbb{E}_0[\Lambda(\varphi^\ast-\psi)]\geq c\mathbb{E}_0(\varphi^\ast-\psi)=c(\alpha-\mathbb{E}_0\psi)\geq0$.\\
\hlt{Structural Commentary}: Null controls size constraint, while alternative supplies power objective.\\
\hlt{Key Insight}: Use the rejection budget where $\Lambda=\frac{dP_1}{dP_0}$ is largest.
\end{proof}
\end{theorem}

\begin{remark} \label{rem:null-alternative-asymmetry} \hlt{Choosing the Null}\\
At level $a$, for simple $P_0$ versus $P_1$, the \hlt{Neyman--Pearson problem} is $\sup\{\mathbb{E}_1\varphi:0\leq\varphi\leq1,\ \mathbb{E}_0\varphi\leq a\}$.\\
Thus $H_0$ is the protected claim: choose it so the false-rejection probability $\mathbb{E}_0\varphi$ is the error to be bounded by $a$.\\
The choice of $H_0$ is substantive, not determined by $\{P_0,P_1\}$.
\end{remark}

\begin{definition} \label{def:uniformly-most-powerful-test} \hlt{Uniformly Most Powerful Test}\\
For a composite alternative $\Theta_1$, a level-$\alpha$ test $\varphi$ is \hlt{uniformly most powerful} if, for every $\theta\in\Theta_1$, it has power at least as large as every other level-$\alpha$ test at that same $\theta$.
\end{definition}

\begin{example} \label{ex:normal-mean-np-one-sided-known-variance} \hlt{Normal Mean, Known Variance}\\
Let $X_1,\ldots,X_n$ be independent $N(\mu,\sigma^2)$ variables with known $\sigma^2$. For $H_0:\mu=\mu_0$ against $H_1:\mu=\mu_1>\mu_0$, the likelihood ratio is increasing in $\overline{X}$, so the level-$\alpha$ test is $\varphi^\ast=\mathbf{1}_{\{\overline{X}>\mu_0+\frac{z_{1-\alpha}\sigma}{\sqrt n}\}}$, where $z_{1-\alpha}$ is the $N(0,1)$ quantile. Since the cutoff is independent of $\mu_1$, $\varphi^\ast$ is uniformly most powerful for $H_1:\mu>\mu_0$; no single likelihood-ratio direction is uniformly most powerful for $H_1:\mu\neq\mu_0$.
\end{example}

\begin{definition} \label{def:generalized-likelihood-ratio-test} \hlt{Generalized Likelihood-Ratio Test}\\
For $H_0:\theta\in\Theta_0\subseteq\Theta$, the \hlt{generalized likelihood ratio} is $\lambda(x)\defeq\frac{\sup\limits_{\theta\in\Theta_0}L_x(\theta)}{\sup\limits_{\theta\in\Theta}L_x(\theta)}\in[0,1]$.\\
A GLR test rejects $H_0$ for small $\lambda(x)$, typically $\lambda(x)\leq c_\alpha$ with $c_\alpha$ chosen to give level $\alpha$.
\end{definition}

\begin{theorem} \label{thm:wilks-likelihood-ratio-chi-square} \hlt{Large-Sample Likelihood-Ratio Calibration}
\begin{enumerate}[label=(\roman*)]
\setlength{\itemsep}{0pt}
\item Let $H_0:\theta\in\Theta_0\subseteq\Theta$, $\dim\Theta=d$, $\dim\Theta_0=d_0$, $r\defeq d-d_0$, and $X^{(n)}\sim P_{\theta_0}$ with $\theta_0\in\Theta_0$.
\item Let $L_n(\theta)\defeq e^{\ell_n(\theta)}$ and $\lambda_n\defeq\frac{\sup\limits_{\theta\in\Theta_0}L_n(\theta)}{\sup\limits_{\theta\in\Theta}L_n(\theta)}$.
\item Assume local coordinates with $\theta_0=0$, $\Theta_0=\{(u,0):u\in\R^{d_0}\}$, $h=(u,v)\in\R^{d_0}\times\R^r$,\\
and $\ell_n(\frac{h}{\sqrt n})-\ell_n(0)=h^\top Z_n-\frac{\norm{h}^2}{2}+o_{P_{\theta_0}}(1)$, where $Z_n=(U_n,V_n)\Rightarrow N_d(0,I_d)$, $U_n\in\R^{d_0}$, $V_n\in\R^r$.
\item
Then under $H_0$, $-2\log\lambda_n\Rightarrow\chi_r^2$.
\end{enumerate}
\begin{proof}
\begin{enumerate}[label=(\roman*)]
\setlength{\itemsep}{0pt}
\item Since $L_n=e^{\ell_n}$, $\log\lambda_n=\sup\limits_{\theta\in\Theta_0}\ell_n(\theta)-\sup\limits_{\theta\in\Theta}\ell_n(\theta)$, hence $-2\log\lambda_n=2\{\sup\limits_{\theta\in\Theta}\ell_n(\theta)-\sup\limits_{\theta\in\Theta_0}\ell_n(\theta)\}$.
\item In the local coordinates, write $\theta=\frac{h}{\sqrt n}$ for the full model and $\theta=\frac{(u,0)}{\sqrt n}$ for the null model.\\
With $Q_n(h)\defeq h^\top Z_n-\frac{\norm{h}^2}{2}$ and $Q_{0,n}(u)\defeq Q_n(u,0)=u^\top U_n-\frac{\norm{u}^2}{2}$, (iii) and cancellation of $\ell_n(0)$ give $-2\log\lambda_n=2\{\sup\limits_{h\in\R^d}Q_n(h)-\sup\limits_{u\in\R^{d_0}}Q_{0,n}(u)\}+o_{P_{\theta_0}}(1)$.
\item For $a\in\R^m$, $x^\top a-\frac{\norm{x}^2}{2}=\frac{\norm{a}^2}{2}-\frac{\norm{x-a}^2}{2}$, so supremum over $x\in\R^m$ is $\frac{\norm{a}^2}{2}$ at $x=a$.\\
Applying this with $a=Z_n$ and $a=U_n$ gives $\sup\limits_{h\in\R^d}Q_n(h)=\frac{\norm{Z_n}^2}{2}$ and $\sup\limits_{u\in\R^{d_0}}Q_{0,n}(u)=\frac{\norm{U_n}^2}{2}$,\\
hence $-2\log\lambda_n=\norm{Z_n}^2-\norm{U_n}^2+o_{P_{\theta_0}}(1)=\norm{V_n}^2+o_{P_{\theta_0}}(1)$.
\item Since $Z_n=(U_n,V_n)\Rightarrow N_d(0,I_d)$, marginal convergence gives $V_n\Rightarrow Y$ with $Y\sim N_r(0,I_r)$.\\
Therefore $\norm{V_n}^2\Rightarrow\norm{Y}^2\sim\chi_r^2$ by continuous mapping, and Slutsky gives $-2\log\lambda_n\Rightarrow\chi_r^2$.
\end{enumerate}
\hlt{Structural Commentary}: The likelihood ratio compares unconstrained fit with null-constrained fit.\\
\hlt{Key Insight}: Wilks' statistic is the squared Gaussian fluctuation in the $r=d-d_0$ constrained directions.
\end{proof}
\end{theorem}

\begin{example} \label{ex:normal-mean-glrt-known-variance} \hlt{Two-Sided Normal Mean GLRT}\\
Let $X_1,\ldots,X_n$ be independent $N(\mu,\sigma^2)$ variables with known $\sigma^2$, and test $H_0:\mu=\mu_0$ against $H_1:\mu\neq\mu_0$.\\
Since $\widehat{\mu}=\overline{X}$, $\lambda(X)=\frac{L_n(\mu_0)}{L_n(\overline{X})}=\exp\{-\frac{n(\overline{X}-\mu_0)^2}{2\sigma^2}\}$, so $-2\log\lambda(X)=\frac{n(\overline{X}-\mu_0)^2}{\sigma^2}$.\\
Under $H_0$, $\frac{\sqrt n(\overline{X}-\mu_0)}{\sigma}\sim N(0,1)$; hence $-2\log\lambda(X)\sim\chi_1^2$,\\
and the level-$\alpha$ GLR rejects iff $\frac{n(\overline{X}-\mu_0)^2}{\sigma^2}>\chi^2_{1,1-\alpha}$, equivalently $\abs{\overline{X}-\mu_0}>\frac{z_{1-\frac{\alpha}{2}}\sigma}{\sqrt n}$.
\end{example}

\begin{example} \label{ex:sign-test-binomial-glrt} \hlt{Sign Test as a Binomial GLRT}\\
Let $X\sim\operatorname{Binomial}(n,p)$ and test $H_0:p=\frac{1}{2}$ against $H_1:p\neq\frac{1}{2}$.\\
For observed $x$, $\widehat p=\frac{x}{n}$ and $\lambda(x)=\frac{L_x(\frac{1}{2})}{L_x(\widehat p)}=\frac{(\frac{1}{2})^n}{[(\frac{x}{n})^x(1-\frac{x}{n})^{n-x}]}$, $0^0=1$ at $x=0,n$.\\
Since $\lambda(x)$ decreases as $\abs{x-\frac{n}{2}}$ increases, a level-$\alpha$ GLR rejects on the largest values of $\abs{X-\frac{n}{2}}$, with cutoff chosen under $X\sim\operatorname{Binomial}(n,\frac{1}{2})$; this is the two-sided sign test.
\end{example}

\begin{definition} \label{def:p-value-and-validity} \hlt{$p$-Value}\\
A \hlt{$p$-value} for $H_0:\theta\in\Theta_0$ is a statistic $p:\mathcal{X}\to[0,1]$ such that $\sup\limits_{\theta\in\Theta_0}\mathbb{P}_\theta\{p(X)\leq u\}\leq u$ for every $u\in[0,1]$.\\
Then $\varphi_\alpha(X)\defeq\mathbf{1}_{\{p(X)\leq\alpha\}}$ has size $\alpha(\varphi_\alpha)\leq\alpha$.
\end{definition}

\begin{proposition} \label{prop:continuous-null-p-value-uniform} \hlt{Continuous Null $p$-Values Are Uniform}\\
Let $H_0:\theta\in\Theta_0$ and let $T:\mathcal{X}\to\R$ satisfy $\mathbb{P}_\theta(T\leq t)=F_0(t)$ for all $\theta\in\Theta_0$, where $F_0$ is continuous.\\
For the right-tail convention, set $p(X)\defeq1-F_0(T(X))$.\\
Then $p(X)\sim\operatorname{Uniform}(0,1)$ under every $\theta\in\Theta_0$, and $\varphi_\alpha(X)\defeq\mathbf{1}_{\{p(X)\leq\alpha\}}$ has size $\alpha$ for $\alpha\in[0,1]$.
\begin{proof} Fix $\theta\in\Theta_0$ and set $U\defeq F_0(T(X))$.\\
By probability integral transform, Proposition \ref{prop:probability-integral-and-inverse-transform}, $U\sim\operatorname{Uniform}(0,1)$ under $\mathbb{P}_\theta$,\\
hence $p(X)=1-U\sim\operatorname{Uniform}(0,1)$. Thus $\mathbb{E}_\theta\varphi_\alpha=\mathbb{P}_\theta(p(X)\leq\alpha)=\alpha$ $\forall \theta\in\Theta_0$, so $\alpha(\varphi_\alpha)=\sup\limits_{\theta\in\Theta_0}\mathbb{E}_\theta\varphi_\alpha=\alpha$.\\
\hlt{Structural Commentary}: A continuous null CDF converts the statistic to a canonical uniform scale.\\
\hlt{Key Insight}: Exact calibration is the probability integral transform applied under $H_0$.
\end{proof}
\end{proposition}

\begin{remark} \label{rem:p-values-are-not-posterior-probabilities} \hlt{$p$-Values Are Tail Probabilities}\\
A $p$-value is a null tail probability, e.g. $p(x)=\mathbb{P}_{H_0}\{T(X)\geq T(x)\}$ for a right-tail statistic $T$.\\
It is not $\mathbb{P}(H_0\mid x)$; posterior probabilities require priors and an alternative likelihood, as in Proposition \ref{prop:posterior-odds-likelihood-ratio}.
\end{remark}

\begin{remark} \label{rem:error-control-families} \hlt{Error Rates}\\
For one test $\varphi$, size is $\alpha(\varphi)=\sup\limits_{\theta\in\Theta_0}\mathbb{E}_\theta\varphi$.
For a family of tests, let $R$ be the number of rejections and $V$ the number of false rejections; common targets are $\operatorname{FWER}=\mathbb{P}\{V\geq1\}$ and $\operatorname{FDR}=\mathbb{E}[\frac{V}{R\vee1}]$, while directional or selective rates replace $V$ by the corresponding error count.
\end{remark}
